\documentclass[11pt]{article}
\usepackage[T1]{fontenc}
\usepackage{lmodern}
\usepackage[a4paper,margin=25mm]{geometry}
\usepackage{amsmath,amssymb,amsthm}
\usepackage[hypertexnames=false,hidelinks]{hyperref}
\newcommand{\Z}{\mathbb{Z}}
\newcommand{\V}{\mathsf V}
\newcommand{\E}{\mathsf E}
\newcommand{\loc}[1]{\par\smallskip\noindent{\small\textit{Source: #1}}\par\smallskip}
\setlength{\emergencystretch}{2em}
\begin{document}
\begin{center}
{\large Probability on Trees and Networks\\[2pt]Exercise 6.74 and its Author Comment}\\[4pt]
Russell Lyons\quad Yuval Peres\\[2pt]
Version of 19 August 2026---corrected and updated paperback edition
\end{center}

% BEGIN CONTEXT 6.74
\noindent\textbf{Source and scope (editorial).}
This private reader reproduces selected English author text from the fixed 716-page paperback PDF (21,010,205 bytes). Exercise 6.74 is on physical p.\ 251 (printed p.\ 234), in the collected in-text exercises of Chapter 6. Its complete Comment is on physical p.\ 637 (printed p.\ 620). The edge boundary $\partial_{\E}$ that the exercise quantifies over is defined in the body on physical p.\ 194 (printed p.\ 177). The product inequality that the Comment invokes, and the theorem it names, are printed together on physical p.\ 220 (printed p.\ 203): Theorem 6.22 (Discrete Isoperimetric Inequality), Lemma 6.23 (Discrete Loomis and Whitney Inequality), the projection maps $P_i$, and the sentence tying the two together. All of this is reprinted below as uncounted context.

The separate errata sheet dated 19 August 2026 contains exactly one Chapter 6 entry, the p.\ 616 correction to Exercise 6.34 (``$v$'' to ``$x$'' and the final ``for some $x\in\partial\V_K$''), and no entry for Exercise 6.74 or for its Comment. Accordingly no revision is applied anywhere in this reader, and no printed text here is altered. No source mathematics is repaired or independently certified.

\medskip\noindent\textbf{Literal source features kept as printed (no synchronization, no repair).}
\begin{itemize}
\item The exercise (physical p.\ 251) and the Comment (physical p.\ 637) both print the hypothesis as $|A|<n^{d}/2$. On both pages the exponent contains only $d$: the $d$ is set at script size and raised above the line, while the solidus and the $2$ are set at body size on the baseline (contrast $n^{1/3}$ on the same p.\ 251, where the $1$, the solidus and the $3$ are all at script size). Neither page prints $n^{d/2}$. Both strings are transcribed exactly as printed and neither is altered.
\item The Comment writes the product bound as following from ``Theorem 6.22''. The statement printed under that name on physical p.\ 220 is the edge-isoperimetric inequality $|\partial_{\E}A|\ge 2d\,|A|^{(d-1)/d}$; the product inequality $\prod_j|P_j(A)|\le|P_m(A)|^{d}$ displayed in the Comment is Lemma 6.23's, which is printed on the same page and from which Theorem 6.22 is derived there. The printed attribution is kept verbatim.
\item The exercise says that $A$ is ``a subgraph of the box'', while the Comment writes $A\subset\{0,\dots ,n-1\}^{d}$. Both are kept as printed; no reading is imposed here.
\item The constant $C_d$ of the exercise is not identified with the Comment's explicit $(1-2^{-1/d})$, and no small-$n$ adjustment is inserted.
\end{itemize}

\medskip\noindent\textbf{Use of the exercise in the body of the work.}
No body theorem or proposition of the fixed version cites Exercise 6.74 by number. Theorem 6.22, which the Comment names, is proved in the body beginning on physical p.\ 220 from Lemma 6.23 and is not a special case of the exercise; the exercise states this itself.

\subsection*{Definition of the edge boundary}
\loc{Physical p.\ 194 (printed p.\ 177).}
For $K\subset\V$, we define the edge boundary $\partial_{\E}K$ to be the set of (unoriented) edges that connect $K$ to its complement. Write
\[
|K|_D:=\sum_{x\in K}D(x)
\qquad\text{and}\qquad
|\partial_{\E}K|_c:=\sum_{e\in\partial_{\E}K}c(e).
\]

\subsection*{Theorem 6.22, Lemma 6.23 and the projections: the named input}
\loc{Physical p.\ 220 (printed p.\ 203).}
\textbf{Theorem 6.22.} (Discrete Isoperimetric Inequality) \textit{If $A\subset\Z^d$ is a finite set, then}
\[
|\partial_{\E}A|\ge 2d\,|A|^{\frac{d-1}{d}}.
\]
Observe that the $2d$ constant in the inequality is the best possible, as the example of the $d$-dimensional cube shows: if $A=[0,n)^d\cap\Z^d$, then $|A|=n^d$ and $|\partial_{\E}A|=2dn^{d-1}$. The same inequality without the sharp constant follows from Theorem 6.29.

To prove this inequality, we will develop other very useful tools concerning entropy. For every $1\le i\le d$, define the projection $P_i:\Z^d\to\Z^{d-1}$ simply as the function dropping the $i$th coordinate, that is, $P_i(x_1,\dots ,x_d)=(x_1,\dots ,x_{i-1},x_{i+1},\dots ,x_d)$. Theorem 6.22 will follow easily from the following beautiful inequality of Loomis and Whitney (1949).

\medskip\noindent\textbf{Lemma 6.23.} (Discrete Loomis and Whitney Inequality) \textit{For any finite $A\subset\Z^d$,}
\[
|A|^{d-1}\le\prod_{i=1}^{d}|P_i(A)|.
\]
% END CONTEXT 6.74

\clearpage
% BEGIN PRIMARY STATEMENT 6.74
\section*{Exercise 6.74}
\loc{Physical p.\ 251 (printed p.\ 234).}
\begin{sloppypar}
\noindent\textbf{Exercise 6.74.} Show that there exists a constant $C_d>0$ such that if $A$ is a subgraph of the box $\{0,\dots ,n-1\}^{d}$ and $|A|<n^{d}/2$, then
\[
|\partial_{\E}A|\ge C_d\,|A|^{\frac{d-1}{d}}.
\]
Here, $\partial_{\E}A$ refers to the edge boundary within the box, so this is not a special case of Theorem 6.22.
\end{sloppypar}
% END PRIMARY STATEMENT 6.74

% BEGIN PRIMARY PROOF 6.74
\section*{Complete author Comment on Exercise 6.74}
\loc{Physical p.\ 637 (printed p.\ 620).}
\textbf{6.74.} Let $A\subset\{0,\dots ,n-1\}^{d}$ with $|A|<n^{d}/2$. Let $m$ be such that $|P_m(A)|$ is maximal over all projections, and let
\[
F=\bigl\{a\in P_m(A)\,;\,|P_m^{-1}(a)|=n\bigr\}.
\]
Notice that for any $a\notin F$, there is at least one edge in $\partial_{\E}A$, and for different $a$'s we get disjoint edges, so that $|\partial_{\E}A|\ge|P_m(A)\setminus F|$. By Theorem 6.22, we get
\[
|A|^{d-1}\le\prod_{j=1}^{d}|P_j(A)|\le|P_m(A)|^{d},
\]
and so $|P_m(A)|\ge|A|/|A|^{1/d}\ge 2^{1/d}|A|/n\ge 2^{1/d}|F|$, which together yield
\[
|\partial_{\E}A|\ge|P_m(A)\setminus F|\ge(1-2^{-1/d})|P_m(A)|\ge(1-2^{-1/d})|A|^{\frac{d-1}{d}}.
\]
% END PRIMARY PROOF 6.74

\medskip\noindent\textbf{Cached bibliographic attribution (not an additional source proof).}
The Comment on Exercise 6.74 carries no external bibliographic attribution and names no external work. Its single named input, ``Theorem 6.22'', is internal to this book, printed on physical p.\ 220 together with Lemma 6.23 and both reprinted above as uncounted context.
\end{document}
